\subsection{TENSOR PRODUCT OF TWO LINEAR MAPPINGS}

Let A be a ring, E, E' two right A-modules, F, F' two left A-modules and \(u: E \rightarrow E'\) and \(v: F \rightarrow F'\) two A-linear mappings. It is easily verified that the mapping

\[
(x, y) \mapsto u (x) \otimes v (y)
\]

of \(\mathbf{E} \times \mathbf{F}\) into \(\mathbf{E}' \otimes_{\mathbf{A}} \mathbf{F}'\) is \(\mathbf{Z}\) -bilinear and satisfies conditions (1) of no. 1. By Proposition 1 of no. 1 there thus exists one and only one \(\mathbf{Z}\) -linear mapping \(w: \mathbf{E} \otimes_{\mathbf{A}} \mathbf{F} \to \mathbf{E}' \otimes_{\mathbf{A}} \mathbf{F}'\) such that

\[
w (x \otimes y) = u (x) \otimes v (y)\tag{3}
\]

for \(x \in \mathbf{E}\) , \(y \in \mathbf{F}\) . This mapping is denoted by \(u \otimes v\) (when no confusion can arise) and is called the tensor product of the linear mappings \(u\) and \(v\) .

It follows immediately from (3) that \((u, v) \mapsto u \otimes v\) is a Z-bilinear mapping called canonical

\[
\operatorname{Hom} _ {\mathbf {A}} (\mathrm{E}, \mathrm{E} ^ {\prime}) \times \operatorname{Hom} _ {\mathbf {A}} (\mathrm{F}, \mathrm{F} ^ {\prime}) \rightarrow \operatorname{Hom} _ {\mathbf {Z}} (\mathrm{E} \otimes_ {\mathbf {A}} \mathrm{F}, \mathrm{E} ^ {\prime} \otimes_ {\mathbf {A}} \mathrm{F} ^ {\prime}).
\]

There corresponds to it by Proposition 1 of no. 1 a \(\mathbf{Z}\) -linear mapping called canonical

\[
\operatorname{Hom} _ {\mathbf {A}} (\mathrm{E}, \mathrm{E} ^ {\prime}) \otimes_ {\mathbf {Z}} \operatorname{Hom} _ {\mathbf {A}} (\mathrm{F}, \mathrm{F} ^ {\prime}) \rightarrow \operatorname{Hom} _ {\mathbf {Z}} (\mathrm{E} \otimes_ {\mathbf {A}} \mathrm{F}, \mathrm{E} ^ {\prime} \otimes_ {\mathbf {A}} \mathrm{F} ^ {\prime})\tag{4}
\]

which associates with every element \(u \otimes v\) of the tensor product the linear mapping \(u \otimes v: \mathbf{E} \otimes_{\mathbf{A}} \mathbf{F} \to \mathbf{E}' \otimes_{\mathbf{A}} \mathbf{F}'\) . Note that the canonical mapping (4) is not necessarily injective nor surjective. The notation \(u \otimes v\) can therefore lead to confusion and it will be necessary for the context to indicate whether it denotes an element of the tensor product or a linear mapping.

Further, let \(\mathbf{E}''\) be a right A-module, \(\mathbf{F}''\) a left A-module and \(u':\mathbf{E}'\to\mathbf{E}''\) , \(v':\mathbf{F}'\to\mathbf{F}''\) A-linear mappings; it follows from (3) that

\[
\left(u ^ {\prime} \circ u\right) \otimes \left(v ^ {\prime} \circ v\right) = \left(u ^ {\prime} \otimes v ^ {\prime}\right) \circ (u \otimes v).\tag{5}
\]
% semantic-fragments: legacy-input-seam
\relax{}
